\section{Simulation algorithms}\label{supp:algorithm}

\begin{algorithm}[!ht]
\caption{Sampling a prevalence survey of one exposure group from the stationary distribution of background result B1 (comparison of the estimators, Aim 4).}
\label{alg:sim}
\begin{algorithmic}[1]
\Require expected group population size $n$; incidence density $I$; emigration rate $\gamma$; duration distribution $F$ with mean $\bar D$
\Ensure healthy count $H$; prevalent count $P$; times since onset $A_1,\dots,A_P$
\State $\bar H\gets n/(1+I\bar D)$; $\bar P\gets I\bar D\bar H$
\State Draw independently $H\sim\mathrm{Poisson}(\bar H)$ and $P\sim\mathrm{Poisson}(\bar P)$
\For{$k=1,\dots,P$}
  \State Draw the backward time $A_k$ from the density $g(a)=S(a)/\bar D$, $a\ge0$
\EndFor
\State \Return $H$, $P$ and the case-specific times since onset $A_k$
\end{algorithmic}
\end{algorithm}

Algorithm~\ref{alg:sim} was used for the comparison of the estimators (scenario sets A to D). In scenario set C, the backward times of the exposed cases were drawn from the density proportional to $S_0(a)e^{-\beta a}$. The algorithm samples the stationary distribution described in background result B1, so this algorithm alone cannot verify that sampling identity. Aims 1 to 3, the checks of the two propositions and the scenarios with a recent start, immigration of prevalent cases, growth and changing incidence rates used simulations of individual histories (Algorithm~\ref{alg:orig} and its variants), in which no count or time since onset is drawn from a derived distribution.

\begin{algorithm}[!ht]
\caption{Simulation of the individual histories of one dynamic population (Aims 1 to 3).}
\label{alg:orig}
\begin{algorithmic}[1]
\Require $\alpha$, $\gamma$, $N_0$, $t_{\max}$, $I$, $\bar D$, $dstr$
\Statex \Comment{$\alpha$, per capita rate of immigration; $\gamma$, per capita rate of emigration; $N_0$, size of the initial population; $t_{\max}$, maximum follow-up time; $I$, disease incidence density rate; $\bar D$, average duration of the disease; $dstr$, distribution used to generate the duration time}
\Ensure the population dynamic until time $t_{\max}$: a record of the disease time and status of the individuals who have been part of the population at any point during the study period
\State \textbf{Rates calculation:}
\State $prv\gets I\bar D/(1+I\bar D)$ \Comment{prevalence calculated from the equation $P/H=I\times\bar D$}
\State $H\gets N_0(1-prv)$ \Comment{expected size of the healthy pool}
\State $a\gets H\times\alpha$ \Comment{rate of immigration to the healthy population}
\State \textbf{Simulating the dynamics of the population:}
\State the $N_0$ subjects of the initial population enter at $t_0=0$
\State $t\gets 0$
\Repeat
  \State generate a waiting time $t_{next}$ from $\exp(a)$;\quad $t\gets t+t_{next}$
  \State add a new subject with $t_0=t$ \Comment{immigration}
\Until{$t\ge t_{\max}$}
\For{each subject}
  \State generate a potential emigration time $t_{emig}\gets t_0+\exp(\gamma)$ \Comment{emigration}
  \State generate a potential disease incidence time $td_{start}\gets t_0+\exp(I)$
  \If{$td_{start}<t_{emig}$}
     \State $D\gets 1$
     \State generate the duration time $du$ from $dstr$: $du\gets\bar D$ (constant), $du\sim\mathrm{unif}(\bar D-3,\bar D+3)$ (uniform) or $du\sim\exp(1/\bar D)$ (exponential)
     \State $td_{end}\gets td_{start}+du$;\quad $t_{end}\gets td_{end}$
  \Else
     \State $D\gets 0$;\quad $t_{end}\gets t_{emig}$ \Comment{the subject leaves the population before the disease starts}
  \EndIf
\EndFor
\State \textbf{Prevalence study on day $t$:} $H(t)\gets$ number of subjects with $t_0\le t<\min(td_{start},t_{emig})$;\quad $P(t)\gets$ number of subjects with $D=1$ and $td_{start}\le t<td_{end}$
\end{algorithmic}
\end{algorithm}

Algorithm~\ref{alg:orig} is the algorithm of the original study with two changes. In the original algorithm emigration events occurred at the constant total rate $H\times\gamma$ and removed a healthy subject chosen at random; here each healthy subject emigrates at the per capita rate $\gamma$, which is the model of the main text. The original algorithm set the end of the disease to the larger of $td_{start}+du$ and $t_{\max}$; the end of the disease is $td_{start}+du$. For Aims 1 to 3 each replicate is one run of the algorithm with a prevalence study on day 2000 (script \path{python_original_reproduction.py}). In the variants used for Aim 4 there is no initial cohort: healthy individuals enter as a Poisson process over a period before the survey that is long relative to the durations and to the time spent at risk, and the incidence or entry rates may depend on time (script \path{python_individual_histories.py}).


